\originalchapter{方格Ising的精确自由能}
\label{ch:ising2}
\source{8.334 L18--L19给Kac--Ward精确自由能并说明费米传递解释；以下Jordan--Wigner对角化为编者补足的等价计算链。}
\originalsection{行传递与对偶耦合}
考虑无限二维方格、零场、各向同性最近邻$J>0$。一行$N$自旋间的垂直键矩阵每站为$\left(\begin{smallmatrix}\e^K&\e^{-K}\\\e^{-K}&\e^K\end{smallmatrix}\right)=\sqrt{2\sinh2K}\,\e^{K^*\sigma^x}$，$\tanh K^*=\e^{-2K}$。水平键对角，因此对称行传递矩阵可写
\begin{equation}
V=(2\sinh2K)^{N/2}\e^{(K/2)\sum\sigma_j^z\sigma_{j+1}^z}\e^{K^*\sum\sigma_j^x}\e^{(K/2)\sum\sigma_j^z\sigma_{j+1}^z}.
\end{equation}
每个因子由显式矩阵恒等式得到，$Z=\Tr V^{N_y}$，大$N_y$时最大本征值控制自由能。
\originalsection{Majorana化与每模对角化}
为了不把“化为自由费米子”当黑箱，取一行的串算符$S_j=\prod_{\ell<j}\sigma_\ell^x$，定义两个Hermitian Majorana算符
\begin{equation}
a_j=S_j\sigma_j^z,\qquad b_j=S_j\sigma_j^y,
\qquad\{a_j,a_k\}=\{b_j,b_k\}=2\delta_{jk},\quad\{a_j,b_k\}=0.
\end{equation}
沿串移动一个站点会得到一次反对易；同站用Pauli乘法直接核对$\sigma_j^x=\ii a_jb_j$及$\sigma_j^z\sigma_{j+1}^z=\ii b_ja_{j+1}$。故上节$V=PQP$的两种指数均为Majorana双线性，其中
\begin{equation}
P=\exp\!\left(\frac{\ii K}{2}\sum_jb_ja_{j+1}\right),\qquad
Q=\exp\!\left(\ii K^*\sum_ja_jb_j\right).
\end{equation}
周期自旋行的末键带总宇称$\mathcal P=\prod_j\sigma_j^x$，在每个固定宇称扇区选对应的周期或反周期Majorana边界。以下$q$为该扇区允许的动量；两个动量网格相差半格，体极限相同。

取$a_q=N^{-1/2}\sum_j\e^{-\ii qj}a_j$，$b_q$同。由$[\ii\gamma_1\gamma_2,\gamma_1]=-2\ii\gamma_2$与另一条对易式，可逐项计算$P,Q$共轭作用于列向量$(a_q,b_q)^{\mathsf T}$。记$c=\cosh K,s=\sinh K$以及$c_*=\cosh2K^*,s_*=\sinh2K^*$，两种作用矩阵分别为
\begin{equation}
H(q)=\begin{pmatrix}c&\ii s\e^{-\ii q}\\-\ii s\e^{\ii q}&c\end{pmatrix},
\qquad
Q_*(q)=\begin{pmatrix}c_*&-\ii s_*\\\ii s_*&c_*\end{pmatrix}.
\end{equation}
例如$Pa_jP^{-1}=ca_j+\ii s b_{j-1}$，$Pb_jP^{-1}=cb_j-\ii s a_{j+1}$；Fourier变换正给$H(q)$中的两个相位。对称行传递的$2\times2$块为$M(q)=H(q)Q_*(q)H(q)$。两因子行列式均为1，因此$\det M=1$；利用迹的循环性，$\operatorname{tr}M=\operatorname{tr}[Q_*H^2]$，且$H^2$只需把$K$改成$2K$。乘开两项对角元，得到
\begin{equation}
\frac12\operatorname{tr}M(q)=\cosh\gamma_q
=\cosh2K\cosh2K^*-\sinh2K\sinh2K^*\cos q.
\end{equation}
因$\det M=1$，其本征值必为$\e^{\pm\gamma_q}$；对偶关系$\sinh2K\sinh2K^*=1$又把上式改写为$\cosh\gamma_q=\cosh2K\coth2K-\cos q$。每个正能复费米模的两态传递本征值相差$\e^{\gamma_q}$，最大本征态在每模取$+\gamma_q/2$，故
\begin{equation}
\log\lambda_{\max}=\frac N2\log(2\sinh2K)+\frac12\sum_q\gamma_q.
\end{equation}
有限行的宇称投影会调整允许动量和个别态，却不改上述每站体自由能；没有投影时不可把两个完整费米Fock空间重复计为原自旋空间。

因此每站精确自由能为
\begin{equation}
-\beta f=\frac12\log(2\sinh2K)+\frac1{2\pi}\int_0^\pi\operatorname{arcosh}[\cosh2K\coth2K-\cos q]\dd q.
\end{equation}
高温$K\to0$两个对数奇异相消，给$\log2$；低温给$2K$，即每站两条满意键的基态能$-2J$。两个边界极限是公式归一的直接检查。用$t=\tanh K$将此一维$\operatorname{arcosh}$积分还原为第\ref{ch:duality}章的双动量行列式积分，两式必须给相同自由能；例如$K=0.2,0.4,0.7$时两式数值差分别小于$3\times10^{-15}$。
\originalsection{能隙、相关长度与热容}
$\gamma_0=2|K-K^*|$，只在$K=K^*$闭合，故精确临界满足$\sinh2K_c=1$。近临界，$\gamma_q\simeq\sqrt{m^2+q^2}$，$m=2|K-K^*|\propto|T-T_c|$，因此$\xi\sim m^{-1}$、$\nu=1$。奇异自由能来自$\int_0^\Lambda\sqrt{m^2+q^2}\dd q$的$-m^2\log|m|$项，二阶温度导数给对数热容，$\alpha=0$带对数，而非有限平均场跳跃。

零场自发磁化的精确关联结果为$m(T)=[1-\sinh^{-4}(2K)]^{1/8}$于$K>K_c$，高温0，故$\beta=1/8$。这一磁化公式需要非局域自旋相关的Toeplitz/Pfaffian分析；本章上述单粒子传递谱完整求出自由能，却没有把它冒称为磁化公式的全证明。课程L19引用该标准结果，本稿保留这个证明范围。

临界自旋相关$G(r)\sim r^{-1/4}$给$\eta=1/4$，配尺度关系得$\gamma=7/4,\delta=15$。这些不同于近四维一回路外推，说明二维强涨落固定点需要真正的非Gaussian结构。含外场、非最近邻或一般三维Ising没有此自由费米传递解，不能把零场精确公式用于任意Ising模型。
\originalsection{可计算的数值核对}
随卷计算用稳定的$\operatorname{arcosh}$与数值积分对照高低温展开，并以小周期方格枚举所有$2^{L^2}$构型核对有限尺寸配分与能量；无限体公式与小$L$有有限尺寸误差，不强求相等。积分在临界端点的导数奇异需单独处理，普通粗网格微分会低估对数峰。
